\documentclass[11pt,a4paper]{article}

\usepackage[T1]{fontenc}
\usepackage{lmodern}
\usepackage[margin=2.5cm]{geometry}
\usepackage{amsmath}
\usepackage{booktabs}
\usepackage{tabularx}
\usepackage{graphicx}
\usepackage{subcaption}
\usepackage[round]{natbib}
\usepackage[hidelinks]{hyperref}

\graphicspath{{figures/}}
\setlength{\emergencystretch}{1.5em}

\usepackage{url}
\DeclareUrlCommand\file{\urlstyle{tt}}  % file names, breakable at / and _
\newcolumntype{L}{>{\raggedright\arraybackslash}X}

\title{Matsuoka--Nakai UMAT --- Test Report}
\author{ConstitutiveModels}
\date{October 5, 2026}

\begin{document}
\maketitle

% ---------------------------------------------------------------------------
\section{Scope}

This report summarises the review and verification of the Matsuoka--Nakai UMAT
\file{c_models/matsuoka_nakai/matsuoka_nakai.c} and its yield surface module
\file{c_models/yield_surfaces/matsuoka_nakai_surface.c}. The unit tests are in
\file{tests/test_matsuoka_nakai.py}; the figures are produced by
\file{docs/mn_test_report/make_figures.py}, which drives the UMAT along element test paths and
controls stresses by Newton iteration on the strain increment where needed. The library under test
is \file{build_C/lib/matsuoka_nakai.dll}, built from the current sources with
\file{Makefile.windows}.

The model is linear elastic, perfectly plastic:
\begin{itemize}
  \item isotropic linear elasticity with $E$ and $\nu$;
  \item the yield criterion of \citet{Matsuoka_1974}, $I_1 I_2 / I_3 = 9 + 8\tan^2\varphi$, for the
        principal stresses shifted by $a = c\cot\varphi$, in the formulation of
        \citet{Lagioia_2016}:
        \begin{equation}
          f = M p - K + J\,\alpha \cos\left(\tfrac{1}{3}\arccos(\beta\sin 3\theta)\right)
          \label{eq:f}
        \end{equation}
        with $p$ the mean stress (tension positive), $J = \sqrt{J_2}$,
        $\sin 3\theta = \tfrac{3\sqrt{3}}{2} J_3 / J_2^{3/2}$ and, with $s = \sin\varphi$,
        \begin{equation}
          M = \frac{2\sqrt{3}\,s}{3 - s},\quad
          K = \frac{2\sqrt{3}\,c\cos\varphi}{3 - s} = M a,\quad
          \alpha = \frac{2\sqrt{3 + s^2}}{3 - s},\quad
          \beta = \frac{s\,(9 - s^2)}{(3 + s^2)^{3/2}};
          \label{eq:constants}
        \end{equation}
  \item non-associated flow, with the plastic potential $g$ of the same form with the dilation
        angle $\psi$ in place of $\varphi$;
  \item an implicit (backward Euler) return mapping by Newton iteration on the stress and the
        plastic multiplier, with the consistent tangent \citep{Simo_1985, deSouzaNeto_2008};
  \item a return to the apex of the cone ($p = -a$) when that is the solution of the return
        mapping, and sub-stepping if the Newton iteration fails (Sec.~\ref{sec:apex}).
\end{itemize}
The surface passes through all six corners of the Mohr--Coulomb pyramid with the same $c$ and
$\varphi$; between the corners it lies outside the pyramid.

The UMAT interface is tension positive (Abaqus convention). In this report, as in the unit tests,
stresses and strains are compression positive, $\sigma_1 \ge \sigma_2 \ge \sigma_3$,
$p = \frac{1}{3}\operatorname{tr}\boldsymbol{\sigma}$, $q = \sqrt{3 J_2}$,
$b = (\sigma_2 - \sigma_3)/(\sigma_1 - \sigma_3)$ and the mobilised friction angle is
$\sin\varphi_m = (\sigma_1 - \sigma_3)/(\sigma_1 + \sigma_3 + 2a)$. All tests use
$E = 20$~MPa, $\nu = 0.3$, $c = 5$~kPa ($a = 8.66$~kPa), $\varphi = 30^\circ$ and
$\psi = 10^\circ$, unless stated otherwise.

% ---------------------------------------------------------------------------
\section{Review of the implementation}
\label{sec:review}

The review covers the version of commit \texttt{96f5efc} (branch \texttt{update\_matsuoka\_nakai}).
Table~\ref{tab:findings} lists the problems that were found and how they were corrected.
Table~\ref{tab:before-after} compares that version with the corrected one; the values are printed
by \texttt{make\_figures.py -{}-compare <library>}.

\begin{table}[h]
  \centering
  \caption{Findings of the review and corrections.}
  \label{tab:findings}
  \small
  \begin{tabularx}{\textwidth}{@{}p{0.5cm}LL@{}}
    \toprule
    & Finding & Correction \\
    \midrule
    1 & $J_3 = \det\mathbf{s}$ in \file{utils.c} used $s_{xy}^2$ instead of $s_{yz}^2$ in the first
        cofactor. The Lode angle, and therefore the yield function and its gradient, were wrong
        for stress states with shear components, so the response depended on the orientation of
        the axes.
      & Corrected; the same correction is already on branch \texttt{detached4}. \\
    2 & The local iteration of the return mapping updated the plastic multiplier with gradients
        of different iterates (the invariant derivatives of the previous iterate), had no line
        search and an absolute tolerance of $10^{-12}$ on $f$. It converges for small overshoots,
        but diverges for large trial stresses, to stresses of the order of $10^5$~kPa.
      & Newton iteration on $(\boldsymbol{\sigma}, \Delta\gamma)$ with the second derivative of
        $g$, a line search, a relative tolerance of $10^{-10}$, and sub-stepping as fallback. \\
    3 & DDSDDE was $\mathbf{D}^e - \mathbf{D}^e\mathbf{n}_g\,(\mathbf{D}^e\mathbf{n}_f)^T / d$ with
        the denominator $d$ of the trial stress (a second \texttt{denom} declared inside the loop
        hides the first; compiler warning C4456) and gradients of different iterates. It is not
        the tangent of the stress update, so a global Newton iteration converges linearly.
      & Consistent tangent of the implicit return mapping. \\
    4 & No treatment of the apex: trial stresses in tension beyond the apex gave non-finite or
        diverging stresses.
      & Return to the apex $p = -a$ when that solves the return mapping (Sec.~\ref{sec:apex}),
        with $10^{-2}\,\mathbf{D}^e$ as tangent (as for the Hardening Soil model);
        \texttt{STATEV(1)} = 2. \\
    5 & For $\varphi = 0$ the constants of the Tresca criterion were used with $M = K = 0$, so the
        cohesion dropped out of $f$: zero strength. For $\psi = 0$ the plastic potential was a
        smoothed Tresca surface ($\beta = 0.9999$), whereas the Matsuoka--Nakai surface tends to
        the von Mises cylinder for $\psi \to 0$: the flow rule jumped at $\psi = 0$.
      & Constants in closed form, Eq.~\eqref{eq:constants}, continuous in the angle. For
        $\varphi = 0$ the surface is the von Mises cylinder through the corners of the Tresca
        hexagon ($q = 2c$); for $\psi = 0$ the flow is deviatoric and normal to that cylinder. \\
    6 & For $\psi < 0$ the formula $\alpha = \frac{2}{\sqrt{3}}\sqrt{A_1}\,M$ gave $\alpha < 0$,
        which reverses the deviatoric plastic strain.
      & $\alpha > 0$ from Eq.~\eqref{eq:constants}; return mapping for contractant flow
        (Sec.~\ref{sec:apex}). \\
    7 & SSE was incremented by $\boldsymbol{\sigma}_{n+1} : \Delta\boldsymbol{\varepsilon}^e$, twice
        the elastic energy for loading from zero stress.
      & Change of $\frac{1}{2}\boldsymbol{\sigma} : (\mathbf{D}^e)^{-1}\boldsymbol{\sigma}$ (exact). \\
    8 & $\varphi = 90^\circ$ and $\psi = \pm 90^\circ$ were accepted, where $\beta = 1$ and the
        derivatives are singular.
      & $0 \le \varphi < 90^\circ$, $-90^\circ < \psi < 90^\circ$. \\
    \bottomrule
  \end{tabularx}
\end{table}

\begin{table}[h]
  \centering
  \caption{Original (commit \texttt{96f5efc}) and corrected implementation. Increments from
           compressive states with $E = 20$~MPa, $c = 5$~kPa, $\varphi = 30^\circ$, $\psi = 10^\circ$.}
  \label{tab:before-after}
  \small
  \begin{tabularx}{\textwidth}{@{}Lcc@{}}
    \toprule
    Check & Original & Corrected \\
    \midrule
    Deviatoric strain increments of 2\,\% in 49 directions from $p = 100$~kPa:
      $\max|I_1 I_2/(I_3 k) - 1|$ & $1.0$ & $9\times10^{-11}$ \\
    \quad directions with $|I_1 I_2/(I_3 k) - 1| > 10^{-6}$ & 43 of 49 & 0 \\
    General plastic increment: $|I_1 I_2/(I_3 k) - 1|$ & $3.1\times10^{-2}$ & $6\times10^{-12}$ \\
    \quad DDSDDE against finite differences, relative difference & $0.10$ & $4\times10^{-10}$ \\
    \quad same increment in a rotated frame, relative difference & $5.2\times10^{-3}$ & $10^{-16}$ \\
    Isotropic extension from 50~kPa by $\varepsilon_v = -3\,\%$ & $|\sigma| \approx 5\times10^7$~kPa
      & apex, $\sigma = -8.66$~kPa \\
    $\varphi = 0$, $c = 20$~kPa: $q$ in triaxial compression ($2c = 40$~kPa) & $0$ & $40.0$~kPa \\
    Global Newton iterations to $10^{-10}$~kPa (Sec.~\ref{sec:numerics}) & $> 40$ & 6 \\
    SSE / $\frac{1}{2}\boldsymbol{\sigma}:\boldsymbol{\varepsilon}$, elastic loading from zero & 2.0 & 1.0 \\
    \bottomrule
  \end{tabularx}
\end{table}

The corrected yield function coincides with the original one for $0 < \varphi < 90^\circ$ (the
closed forms are algebraically equal to the expressions of \citealp{Lagioia_2016}); it differs only
for $\varphi = 0$. The return mapping, the tangent, the apex return and the plastic potential for
$\psi \le 0$ are new. In triaxial compression and extension the original implementation gives the
correct strength and flow, because the stress stays on a meridian, where its iteration converges
and $J_3$ is not affected by finding~1; this is what the two original unit tests check.

% ---------------------------------------------------------------------------
\section{Test results}

All 37 unit tests pass (pytest, October 5, 2026). Of the 35 new tests, 21 fail with the original
implementation (marked $\times$); the two original tests pass with both.

\begin{table}[h]
  \centering
  \footnotesize
  \begin{tabularx}{\textwidth}{@{}>{\raggedright\arraybackslash}p{0.46\textwidth}Lcc@{}}
    \toprule
    Test & Checks & Result & Original \\
    \midrule
    \file{strain_controlled_compression_triaxial, ..._tension_triaxial} & original tests, strength at the meridians & pass & pass \\
    \file{returned_stress_satisfies_matsuoka_nakai_criterion_at_all_lode_angles} & $I_1 I_2/I_3 = 9 + 8\tan^2\varphi$ at all Lode angles & pass & $\times$ \\
    \file{triaxial_strength_equals_mohr_coulomb} ($\times 4$) & compression, extension; $c = 0$, 10~kPa & pass & pass \\
    \file{dilatancy_at_triaxial_failure} ($\times 6$) & $4\sin\psi/(3 \mp \sin\psi)$, $\psi = 0$, 10, $25^\circ$ & pass & pass \\
    \file{dilatancy_at_triaxial_failure} ($\times 2$) & idem, $\psi = -10^\circ$ (contractant) & pass & $\times$ \\
    \file{plane_strain_intermediate_stress_at_failure} ($\times 2$) & $\sigma_2$ at failure, $\psi = \varphi$ and $\psi = 0$ & pass & $\times$ \\
    \file{undrained_triaxial_compression} ($\times 2$) & $p'$ constant for $\psi = 0$; failure line & pass & pass \\
    \file{undrained_contractant_flow_reaches_the_apex} ($\times 2$) & $\psi < 0$: down the failure line to the apex & pass & $\times$ \\
    \file{zero_friction_angle_gives_von_mises} & $q = 2c$ for $\varphi = 0$ & pass & $\times$ \\
    \file{tangent_matches_numerical_derivative} ($\times 6$) & DDSDDE vs.\ finite differences & pass & $\times$ \\
    \file{principal_frame_invariance} & rotated loading gives the rotated stress & pass & $\times$ \\
    \file{tension_beyond_apex_returns_to_apex} ($\times 3$) & apex stress and tangent, elastic reloading & pass & $\times$ \\
    \file{large_increments_return_to_the_surface} ($\times 2$) & 200 random increments, $\psi = 10^\circ$ and $-10^\circ$ & pass & $\times$ \\
    \file{energy_terms} & SSE and SPD & pass & $\times$ \\
    \file{unloading_from_failure_is_elastic} & elastic unloading and tangent & pass & pass \\
    \file{stress_update_is_continuous_in_strain_increment} & no jumps in the stress update & pass & pass \\
    \bottomrule
  \end{tabularx}
\end{table}

The tangent tests compare DDSDDE after one large plastic increment in a general direction (all six
strain components, $|\Delta\boldsymbol{\varepsilon}| \approx 1.8\,\%$) with central differences of
the stress update: the relative difference is below $10^{-6}$ (tolerance of the test) for
$c = 0$ and 5~kPa and $\psi = 0$, 10 and $30^\circ$; it is $4\times10^{-10}$ in
Table~\ref{tab:before-after}. For $\psi \ne \varphi$ the tangent is not symmetric, for
$\psi = \varphi$ it is symmetric to $10^{-12}$.

% ---------------------------------------------------------------------------
\section{Yield surface}

Single deviatoric strain increments of 2\,\% from an isotropic state of 100~kPa, in 49 directions
of the deviatoric plane, are all plastic and all return onto the Matsuoka--Nakai surface:
$|I_1 I_2/(I_3 k) - 1| \le 9.2\times10^{-11}$, with $k = 9 + 8\tan^2\varphi$ and the
shifted stresses (Fig.~\ref{fig:surface}a). Sweeping one sextant for $\varphi = 20$, 30 and
$40^\circ$, the mobilised friction angle at failure follows the closed-form solution of the
criterion to $10^{-9}$ degrees (Fig.~\ref{fig:surface}b). It equals $\varphi$ at $b = 0$ and
$b = 1$ (the Mohr--Coulomb corners) and is largest in between: 22.94, 34.11 and $44.94^\circ$, at
$b = 0.399$, 0.347 and 0.293.

\begin{figure}[h]
  \centering
  \begin{subfigure}[t]{0.48\textwidth}
    \centering
    \includegraphics[width=\textwidth]{deviatoric_section.pdf}
    \caption{deviatoric plane, shifted stresses normalised by $p + a$, $\varphi = 30^\circ$}
  \end{subfigure}\hfill
  \begin{subfigure}[t]{0.48\textwidth}
    \centering
    \includegraphics[width=\textwidth]{strength_b.pdf}
    \caption{mobilised friction angle at failure vs.\ $b$}
  \end{subfigure}
  \caption{Stresses returned by the UMAT from single large deviatoric strain increments, against
           the Matsuoka--Nakai criterion \citep{Matsuoka_1974} and the Mohr--Coulomb criterion.}
  \label{fig:surface}
\end{figure}

For $\varphi = 0$ the surface is the von Mises cylinder $q = 2c$: triaxial compression, triaxial
extension and simple shear all fail at $q = 2c$ to $10^{-10}$, at constant $p$
(\file{zero_friction_angle_gives_von_mises}).

% ---------------------------------------------------------------------------
\section{Drained triaxial compression and extension}

Drained triaxial compression and extension start from 100~kPa, with $\sigma_3 = 100$~kPa in
compression and $\sigma_1 = \sigma_2 = 100$~kPa in extension. At failure,
$(\sigma_1 + a)/(\sigma_3 + a) = N_\varphi = (1 + \sin\varphi)/(1 - \sin\varphi) = 3$ to
$10^{-15}$, so $q_f = 217.32$~kPa in compression and 72.44~kPa in extension, as for
Mohr--Coulomb (Fig.~\ref{fig:triaxial}a). The dilation angle does not affect $q$ in these tests.

At failure the stress is stationary, so the strain increments are plastic. On the compression and
extension meridians the gradient of the plastic potential has no Lode angle component, and
\begin{equation}
  -\frac{\Delta\varepsilon_v}{\Delta(\varepsilon_1 - \varepsilon_3)} =
  \frac{4\sin\psi}{3 - \sin\psi} \quad\text{(compression)}, \qquad
  \frac{4\sin\psi}{3 + \sin\psi} \quad\text{(extension)}.
\end{equation}
This equals the dilatancy of the Mohr--Coulomb potential at its corners with both adjacent planes
equally active \citep{Koiter_1960}. The UMAT follows it to $1.4\times10^{-15}$ for
$\psi = -20$ to $30^\circ$ (Fig.~\ref{fig:triaxial}b--c); for $\psi < 0$ the soil contracts at
failure.

\begin{figure}[h]
  \centering
  \begin{subfigure}[t]{0.48\textwidth}
    \centering
    \includegraphics[width=\textwidth]{triaxial_q.pdf}
    \caption{deviator stress vs.\ shear strain}
  \end{subfigure}\hfill
  \begin{subfigure}[t]{0.48\textwidth}
    \centering
    \includegraphics[width=\textwidth]{triaxial_ev.pdf}
    \caption{volumetric strain, compression}
  \end{subfigure}

  \vspace{1ex}
  \begin{subfigure}[t]{0.48\textwidth}
    \centering
    \includegraphics[width=\textwidth]{triaxial_dilatancy.pdf}
    \caption{dilatancy ratio at failure vs.\ $\psi$}
  \end{subfigure}
  \caption{Drained triaxial compression and extension from 100~kPa.}
  \label{fig:triaxial}
\end{figure}

% ---------------------------------------------------------------------------
\section{Plane strain}

In plane strain ($\varepsilon_y = 0$, $\sigma_z = 100$~kPa, compression along $x$) the stress
reaches the surface at $b = \nu = 0.3$ and then moves along it until the plastic strain in the
plane-strain direction vanishes, $\partial g/\partial\sigma_2 = 0$ (Fig.~\ref{fig:plane-strain}).
The UMAT reaches the stationary state of this condition to $10^{-12}$:
\begin{itemize}
  \item $\psi = \varphi$ (associated): $\sigma_2 + a = \sqrt{(\sigma_1 + a)(\sigma_3 + a)}$, the
        classical result for the Matsuoka--Nakai criterion. The stress ratio then follows from
        $t + 1/t = \sqrt{k} - 1$ with $t^2 = (\sigma_1 + a)/(\sigma_3 + a)$, giving
        $b = 1/(1 + t) = 0.3466$ and $\varphi_m = 34.11^\circ$: the largest strength over $b$
        (Fig.~\ref{fig:surface}b), because $\partial f/\partial\sigma_2 = 0$ there;
  \item $\psi = 0$: the potential is the von Mises cylinder, so $\sigma_2 = (\sigma_1 + \sigma_3)/2$,
        $b = 0.5$ and $\varphi_m = 33.69^\circ$;
  \item $\psi = 10^\circ$: $b = 0.4495$ and $\varphi_m = 33.91^\circ$, from the condition solved
        independently of the UMAT.
\end{itemize}
The plane strain strength is therefore 3.7 to $4.1^\circ$ higher than the triaxial
$\varphi = 30^\circ$, and depends on $\psi$.

\begin{figure}[h]
  \centering
  \begin{subfigure}[t]{0.48\textwidth}
    \centering
    \includegraphics[width=\textwidth]{plane_strain_b.pdf}
    \caption{$b$ vs.\ axial strain}
  \end{subfigure}\hfill
  \begin{subfigure}[t]{0.48\textwidth}
    \centering
    \includegraphics[width=\textwidth]{plane_strain_phi.pdf}
    \caption{mobilised friction angle vs.\ axial strain}
  \end{subfigure}
  \caption{Plane strain compression from 100~kPa (120 steps). Dashed: stationary state at
           failure from $\partial g / \partial\sigma_2 = 0$.}
  \label{fig:plane-strain}
\end{figure}

% ---------------------------------------------------------------------------
\section{Undrained triaxial compression}

Undrained (isochoric) triaxial compression from $p' = 100$~kPa: $\mathrm{d}p' = -K\,\mathrm{d}\varepsilon_v^p$,
so $p'$ changes only through plastic volume change. With $\psi = 0$, $p'$ stays at 100~kPa (to the
last digit printed, $10^{-10}$~kPa) and the path ends on the failure line
$q = \frac{6\sin\varphi}{3 - \sin\varphi}(p' + a) = 1.2\,(p' + a)$ at $q = 130.4$~kPa. With
$\psi = 10^\circ$ the soil dilates at failure and the path climbs along the failure line, to
$p' = 306.5$~kPa at $\varepsilon_1 = 5\,\%$ (Fig.~\ref{fig:undrained}); without a dilatancy
cut-off this continues indefinitely. With $\psi = -5^\circ$ the soil contracts, and the path
descends along the failure line (static liquefaction) until it reaches the apex,
$p' = -a = -8.66$~kPa and $q = 0$, at $\varepsilon_1 = 3.85\,\%$, where it stays. On the failure
line all paths satisfy $q = 1.2\,(p' + a)$ to $2\times10^{-15}$.

\begin{figure}[h]
  \centering
  \begin{subfigure}[t]{0.48\textwidth}
    \centering
    \includegraphics[width=\textwidth]{undrained_pq.pdf}
    \caption{$p'$--$q$}
  \end{subfigure}\hfill
  \begin{subfigure}[t]{0.48\textwidth}
    \centering
    \includegraphics[width=\textwidth]{undrained_q.pdf}
    \caption{$\varepsilon_1$--$q$}
  \end{subfigure}
  \caption{Undrained triaxial compression from $p' = 100$~kPa with contractant
           ($\psi = -5^\circ$), non-dilatant and dilatant flow.}
  \label{fig:undrained}
\end{figure}

% ---------------------------------------------------------------------------
\section{Apex and contractant flow}
\label{sec:apex}

The return to the apex $\boldsymbol{\sigma}_A = -a\,\mathbf{I}$ solves the return mapping when its
plastic strain $(\mathbf{D}^e)^{-1}(\boldsymbol{\sigma}^{tr} - \boldsymbol{\sigma}_A)$ is a flow
direction of the plastic potential at its apex, $\lambda\,(\frac{1}{3}M_g\mathbf{I} + \mathbf{n})$
with $\lambda \ge 0$ and $h^*(\mathbf{n}) \le 1$. Here $h^*$ is the support function of the
deviatoric section of $g$ (the dual norm of its deviatoric term in Eq.~\eqref{eq:f}), which the
UMAT evaluates by a maximisation over the Lode angle; it agrees with a brute-force maximisation
over random deviatoric directions to $10^{-13}$. In terms of the trial stress (tension positive):
\begin{equation}
  \lambda = \frac{p^{tr} - p_A}{K M_g} \ge 0, \qquad h^*(\mathbf{s}^{tr}) \le 2G\lambda.
\end{equation}
\begin{itemize}
  \item $\psi > 0$: this region lies beyond the apex and is complementary to the trial stresses
        that return to the cone, so the decision is exact.
  \item $\psi \le 0$: trial stresses beyond the apex cannot be returned to the cone, because the
        flow does not reduce $p$. They go to the apex, although the volumetric plastic strain then
        does not follow the flow rule; this is the usual treatment of the apex.
  \item $\psi < 0$: the region lies on the compressive side of the apex and overlaps with the
        trial stresses that return to the cone. The return to the cone is used where it exists,
        as the solution that is continuous for small increments. For trial stresses far outside
        the cone in the deviatoric direction neither return exists; for a Drucker--Prager cone
        both conditions reduce to $J^{tr} \le G\alpha_g(p_A - p^{tr})/(K|M_g|)$. This also holds at
        the apex itself under continued shearing. The increment is then divided into sub-steps,
        and in the smallest sub-steps (1/64) such trial stresses are returned to the apex, where
        contractant flow leads anyway.
\end{itemize}
The tension tests and Fig.~\ref{fig:undrained} ($\psi = -5^\circ$) cover these cases.

% ---------------------------------------------------------------------------
\section{Numerical aspects}
\label{sec:numerics}

\paragraph{Global convergence.} The test of Fig.~\ref{fig:numerics}a mimics a finite element
iteration. From a state at failure (eight steps of plane strain loading), $\Delta\varepsilon_x$
and $\Delta\varepsilon_y$ are prescribed, and $\sigma_z = 100$~kPa, $\sigma_{xy} = 5$~kPa and
$\sigma_{yz} = \sigma_{xz} = 0$ are the targets; the other strain components are solved for by
Newton iteration. With DDSDDE the residual falls quadratically,
$5.4\times10^{-2} \to 1.4\times10^{-5} \to 9\times10^{-13}$~kPa, in six iterations. With the
elastic stiffness it is still 3.1~kPa after 15 iterations, and with the tangent of the original
implementation it does not reach $10^{-10}$~kPa within 40 iterations.

\paragraph{Step size.} On the compression and extension meridians the return is radial in the
deviatoric plane, so the triaxial results above do not depend on the step size. On other paths
the implicit integration has a first-order error in the step size. For plane strain with
$\psi = 10^\circ$ up to $\varepsilon_1 = 10\,\%$ (Fig.~\ref{fig:numerics}b), compared with 1000
steps, the largest difference in $b$ is $1.9\times10^{-2}$, $1.2\times10^{-2}$ and
$2.9\times10^{-3}$ for 5, 20 and 100 steps, and in $q$ (relative to $q_f$)
$3.4\times10^{-3}$, $1.8\times10^{-3}$ and $4.3\times10^{-4}$. The stationary state at failure is
the same for every step size.

\paragraph{Robustness.} The unit tests include 200 random increments with strains of up to the
order of 10\,\%, for $\psi = 10^\circ$ and $-10^\circ$. The review used 4000 random increments with
random parameters ($0 \le \varphi < 50^\circ$, $0 \le \psi \le \varphi$, $0 \le c \le 30$~kPa,
$10^3 \le E \le 10^6$~kPa): every plastic state lies on the surface to $9\times10^{-11}$, and the
plastic strain increment is parallel to the gradient of $g$ at the returned stress to
$10^{-10}$. In addition, 40\,000 increments with extreme parameters and strains
($\varphi$ up to $89^\circ$, $-30^\circ \le \psi \le 89^\circ$, $\nu$ up to 0.49, strain increments
up to about 300\,\%) were integrated without a failure. Sub-stepping was needed in 1196 of them,
all but four with $\psi < 0$; in the tests and figures with $\psi \ge 0$ it was never needed. The
stress update is continuous in the strain increment, also across the elastic--plastic transition
(\file{stress_update_is_continuous_in_strain_increment}).

\begin{figure}[h]
  \centering
  \begin{subfigure}[t]{0.48\textwidth}
    \centering
    \includegraphics[width=\textwidth]{global_newton.pdf}
    \caption{global Newton iteration of a mixed-control increment}
  \end{subfigure}\hfill
  \begin{subfigure}[t]{0.48\textwidth}
    \centering
    \includegraphics[width=\textwidth]{step_size.pdf}
    \caption{plane strain, $\psi = 10^\circ$, with 5, 20 and 100 steps}
  \end{subfigure}
  \caption{Convergence with the consistent tangent, and step-size dependence.}
  \label{fig:numerics}
\end{figure}

% ---------------------------------------------------------------------------
\section{Open points}
\label{sec:open}

\begin{enumerate}
  \item \emph{$\varphi = 0$ and $\psi = 0$.} The surface and the plastic potential now become the
        von Mises cylinder for a zero angle, the limit of the Matsuoka--Nakai criterion. The
        original implementation intended a Tresca surface there, as for Mohr--Coulomb. Triaxial
        results are the same for both, but plane strain with $\psi = 0$ now gives $b = 0.5$ at
        failure. A Tresca potential for $\psi = 0$ would make the flow rule discontinuous at
        $\psi = 0$.
  \item \emph{Negative dilation angle.} The Matsuoka--Nakai criterion is even in the angle, so a
        contractant potential is not defined by it. The constants of Eq.~\eqref{eq:constants} are
        continued analytically, which gives $M < 0$ and a deviatoric section mirrored with respect
        to that of $|\psi|$, as for the Mohr--Coulomb potential with $\psi < 0$; with it the
        dilatancy at the meridians is that of the Mohr--Coulomb corners (Fig.~\ref{fig:triaxial}c).
        Contractant perfect plasticity has no solution of the return mapping for some trial
        stresses (Sec.~\ref{sec:apex}); the apex fallback used there is a pragmatic choice.
  \item \emph{Apex.} With $\psi \le 0$, trial stresses beyond the apex go to the apex and the
        plastic strain then has a volumetric part, although $\psi \le 0$. The tangent at the
        apex is $10^{-2}\,\mathbf{D}^e$, not the exact zero. A convex reformulation that treats
        this region consistently is given by \citet{Panteghini_2014}.
  \item \emph{Sub-stepping.} If the Newton iteration fails for a full increment, the increment is
        divided into $2, 4, \ldots, 64$ sub-steps; DDSDDE is then the tangent of the last
        sub-step and the stress update is not continuous where the fallback starts. If all levels
        fail, the stress is not updated and \texttt{PNEWDT} = 0.5 is requested; this did not occur
        in any of the checks above.
\end{enumerate}

\bibliographystyle{plainnat}
\bibliography{../refs}

\end{document}
